\section{Introduction and scope}
\label{sec:intro}

\subsection{The question}

A conventional search fixes a model, a final state and a mass window before looking at data. The
look-elsewhere effect is then a modest correction applied at the end, and nobody worries much
about it. A growing fraction of the resonance program does not work that way. Model-independent
and data-directed scans~\cite{ddp}, machine-learning anomaly searches~\cite{atlas_ml_anomaly} and
general searches over many event classes~\cite{atlas_general_search} all sweep large parts of the
mass space at once, and for those the trials factor is not a footnote. It is the leading term in
what a discovery has to clear.

A prior question therefore needs a careful answer: how many independent looks does the BSM
resonance program actually contain? The number is usually quoted per analysis, if at all, and
never for the field as a whole. This note works it out from the ground up.

\begin{enumerate}[leftmargin=1.6em,itemsep=2pt]
  \item How many independent bump hunts does covering the public BSM model space imply, and what
        local significance does a $5\sigma$ global discovery therefore cost
        (Section~\ref{sec:budget})?
  \item What does the same arithmetic say about a fully combinatorial scan, and which object
        combinations has ATLAS in fact never scanned (Section~\ref{sec:combinatorial})?
  \item Is the resulting trials factor consistent with the excesses ATLAS has actually reported
        (Section~\ref{sec:excess})?
  \item Given a scan, how should candidates be selected: the largest bin, a threshold, or a
        false-discovery-rate rule (Section~\ref{sec:selection})?
  \item Does a two-stage A/B unblinding scheme make the trials factor cheaper, or only more
        defensible (Section~\ref{sec:ab})?
\end{enumerate}

\subsection{The one-paragraph answer}

The BSM resonance space looks unbounded at the level of models, with dozens of theories and
thousands of signal samples across them, and it collapses to \textbf{46 distinct invariant- or
transverse-mass spectra} once organised by where the bump appears, with lepton flavour and
$b$-tag content counted as part of the mass axis. Counting Gross--Vitells
resolution elements~\cite{grossvitells} over each spectrum's published scan window, the whole
program amounts to $\Nsig \approx 3.7\times10^{3}$ independent looks, rising to $6.6\times10^{3}$
when searches are sliced at realistic event-selection granularity. Because the penalty enters as
$\ln \Nsig$, that program costs a local bar of $\Zloc \approx 6.4$ to $6.5$ for a $5\sigma$ global
discovery, and the figure barely moves under any counting choice we varied. \textbf{Breadth is
cheap.} What is expensive is anything that makes $\Nsig$ uncountable, which is the situation a
neural-network-driven scan naturally creates, and the reason the two-stage unblinding discussion
in Section~\ref{sec:ab} matters even though splitting the data cannot reduce the look-elsewhere
effect.

\subsection{What this note does and does not use}

Every number here derives from two data-free ingredients: a curated list of bump observables with
their published scan windows, floors and resolutions (Appendix~\ref{app:windows}), and a map from
public BSM model classes to the spectra they populate. Both live in one module each in the
accompanying repository (Section~\ref{sec:software}), and neither reads a single experiment-internal
input --- no dataset identifiers, no signal-sample inventory, no unpublished analysis
configuration. Whether a given spectrum happens to have centrally produced Monte-Carlo samples
behind it is a question about a collaboration's production program rather than about the search
budget, and it is deliberately outside the scope of this note; a companion ATLAS-internal note
treats it.

One consequence is worth stating early, because it is what keeps the exercise
non-circular. Every spectrum is counted over the mass range a \emph{published} search family
scans, never over a range inferred from which signal samples exist. A search does not become
cheaper because few samples were generated for it.

\subsection{Conventions}

Throughout, one \emph{search} means one scanned mass spectrum, $n_s$ is the number of effective
independent resolution elements it contains, and $\Nsig = \sum_s n_s$. The local significance
required for a $5\sigma$ global discovery follows from
\begin{equation}
  \Zglob^2 \;=\; \Zloc^2 - 2\ln \Nsig
  \qquad\Longrightarrow\qquad
  \Zloc \;=\; \sqrt{25 + 2\ln \Nsig}\,,
  \label{eq:lee}
\end{equation}
the asymptotic Gross--Vitells relation~\cite{grossvitells} in the fixed-resolution-element
approximation. Masses are in GeV unless stated otherwise, and $\sqrt{s} = 13.6\TeV$.

Bibliographic details for the external searches quoted throughout were checked against the arXiv
metadata records and INSPIRE in August 2026.
